\documentclass[11pt]{article}
\usepackage[margin=0.85in]{geometry}
\usepackage{amsmath,amssymb,amsthm,hyperref}
\hypersetup{colorlinks=true,urlcolor=blue,linkcolor=blue}
\newtheorem{theorem}{Theorem}
\newtheorem{lemma}{Lemma}
\title{The optimal logarithmic interpolation bound\\
\large MathDB 372857: solution of Schmidlin's Conjecture 36}
\author{Principia Math In Progress}
\date{12 September 2026}
\begin{document}
\maketitle

\section*{Task 1. Prove the conjectured interpolation bound}
With the source's definitions, \(x\geq0\), \(0\leq\alpha\leq1\), and
\[
 I(x,\alpha)=\widehat\iota(x,\alpha)
 =\log\frac{\Gamma(x+1)(x+1)^\alpha}{\Gamma(x+\alpha+1)}.
\]
The identity with the integral formulation is Theorem 6 of
\cite{Schmidlin}, p.~6. The domain is defined on pp.~3--4 and
Conjecture 36 is on p.~26.

\begin{theorem}\label{main}
For every \(x\geq0\) and \(\alpha\in[0,1]\),
\begin{equation}\label{target}
 I(x,\alpha)<\frac1{8(x+\alpha)+4}.
\end{equation}
Moreover, \(4\) is the largest possible additive constant in a
bound of the form \(I(x,\alpha)\leq[8(x+\alpha)+C]^{-1}\)
on the full domain with a positive denominator.
\end{theorem}

\begin{lemma}\label{center}
For \(u\geq1\), put
\[
 c(u)=\log\Gamma(u)+\tfrac12\log u-\log\Gamma(u+\tfrac12),
 \qquad s=\frac1{(2u+1)^2}.
\]
Then \(c(u)<H(u)\), where
\begin{equation}\label{H}
 H(u)=\frac1{8(u+1)}-\frac1{192(u+1)^3}+\frac1{630(u+1)^5}
      +\frac{s}{2}+\frac{s^2}{4}+\frac{s^3}{6(1-s)}.
\end{equation}
Also \(c(u)=1/(8u)-1/(192u^3)+O(u^{-5})\) as \(u\to\infty\).
\end{lemma}
\begin{proof}
For the Stirling remainder
\[
 r(z)=\log\Gamma(z)-(z-\tfrac12)\log z+z-\tfrac12\log(2\pi),
\]
Binet's integral \cite[Eq.~5.9.10]{DLMF} gives, for every \(z>0\),
\begin{equation}\label{binet}
 \frac1{12z}-\frac1{360z^3}<r(z)
 <\frac1{12z}-\frac1{360z^3}+\frac1{1260z^5}.
\end{equation}
Indeed, integrate
\(t-t^3/3<\arctan t<t-t^3/3+t^5/5\) against the positive
Binet kernel; these inequalities hold for all \(t>0\) by integrating
the corresponding finite geometric identities for \(1/(1+t^2)\).
The moments are
\(\int_0^\infty t^j/(e^{2\pi t}-1)\,dt=j!\zeta(j+1)/(2\pi)^{j+1}\).

The Gamma duplication formula gives \(c(z)=2r(z)-r(2z)\).
Applying opposite bounds from \eqref{binet} therefore yields
\[
 c(z)<\frac1{8z}-\frac1{192z^3}+\frac1{630z^5},
\]
and also gives the stated asymptotic expansion.
The Gamma recurrence gives
\[
 c(u)-c(u+1)=\log\frac{u+1/2}{\sqrt{u(u+1)}}
            =-\tfrac12\log(1-s).
\]
Finally, for \(0<s<1\), the positive logarithmic series implies
\[
 -\tfrac12\log(1-s)
 \leq\frac{s}{2}+\frac{s^2}{4}+\frac{s^3}{6(1-s)}.
\]
Using the preceding bound for \(c(u+1)\) proves \eqref{H}.
\end{proof}

\begin{lemma}\label{psi}
Let \(\psi=\Gamma'/\Gamma\). For every \(w>0\),
\begin{align}
 \frac1{24w^2}-\frac7{960w^4}
 &\leq\psi(w+\tfrac12)-\log w\leq\frac1{24w^2}, \label{psibound}\\
 \psi'(w+\tfrac12)&\geq\frac1w-\frac1{12w^3}.   \label{trigamma}
\end{align}
\end{lemma}
\begin{proof}
We first prove the global elementary bounds
\begin{equation}\label{sinh}
 \frac1t-\frac t6\leq\frac1{\sinh t}
 \leq\frac1t-\frac t6+\frac{7t^3}{360}\qquad(t>0).
\end{equation}
For the left bound, expand
\(t-(1-t^2/6)\sinh t\) in powers of \(t\).
The coefficient of \(t^{2n+1}\), \(n\geq1\), is
\[
 \frac{(2n)(2n+1)/6-1}{(2n+1)!}\geq0.
\]
For the right bound, the coefficients of degrees \(1,3,5\) in
\((1-t^2/6+7t^4/360)\sinh t-t\) vanish, and for \(n\geq3\)
the coefficient of \(t^{2n+1}\) is
\[
 \frac{1+A(7B/360-1/6)}{(2n+1)!}>0,
 \qquad A=(2n)(2n+1),\quad B=(2n-2)(2n-1)\geq20.
\]
Both series converge everywhere, so multiplication by the positive
\(t\sinh t\) proves \eqref{sinh}.

The digamma integral \cite[Eq.~5.9.13]{DLMF}, combined with
Frullani's integral for \(\log((w+1/2)/w)\), gives
\begin{equation}\label{psiint}
 \psi(w+\tfrac12)-\log w
 =\int_0^\infty e^{-wt}
       \left(\frac1t-\frac1{2\sinh(t/2)}\right)\,dt.
\end{equation}
For clarity, the Frullani term adds
\((e^{-wt}-e^{-(w+1/2)t})/t\) to the integrand from
Eq.~5.9.13, which is exactly \eqref{psiint}.
By \eqref{sinh}, the parenthesized kernel lies between
\(t/24-7t^3/5760\) and \(t/24\).
Integration proves \eqref{psibound}.
Also, differentiating the standard digamma integral
\cite[Eq.~5.9.16]{DLMF} gives
\[
 \psi'(w+\tfrac12)
 =\int_0^\infty\frac{t e^{-wt}}{2\sinh(t/2)}\,dt.
\]
The lower bound in \eqref{sinh} makes this at least
\(\int_0^\infty e^{-wt}(1-t^2/24)\,dt\), which is
\eqref{trigamma}. All integrals converge for \(w>0\), and
differentiation is dominated on each compact positive \(w\)-interval.
\end{proof}

\begin{proof}[Proof of Theorem \ref{main}]
Set \(u=x+1\geq1\), \(\delta=\alpha-\tfrac12\in[-\tfrac12,\tfrac12]\),
and
\[
 F_u(\delta)=\log\Gamma(u)+(\tfrac12+\delta)\log u
             -\log\Gamma(u+\tfrac12+\delta)-\frac1{8(u+\delta)}.
\]
All terms are defined because \(u+\delta\geq1/2\). We have
\begin{align}
 F'_u(\delta)&=\log u-\psi(u+\tfrac12+\delta)
                       +\frac1{8(u+\delta)^2}, \label{F1}\\
 F''_u(\delta)&=-\psi'(u+\tfrac12+\delta)
                       -\frac1{4(u+\delta)^3}
 \leq-\frac1{u+\delta}-\frac1{6(u+\delta)^3}<0. \label{F2}
\end{align}
Thus \(F_u\) is strictly concave on the full interval.

By \eqref{psibound}, the slope at zero satisfies
\begin{equation}\label{slope}
 0<\frac1{12u^2}\leq F'_u(0)\leq
 b(u):=\frac1{12u^2}+\frac7{960u^4}.
\end{equation}
Let
\[
 h=\frac1{10},\qquad v=u+h,\qquad
 K(u)=\frac1v+\frac1{6v^3}.
\]
For \(0\leq\delta\leq h\), \eqref{F2} gives \(F''_u(\delta)\leq-K(u)\).
To locate the maximum, observe the exact identity
\begin{equation}\label{J}
 \frac{K(u)}{10}-b(u)=\frac{P(u)}{960u^4(10u+1)^3},
\end{equation}
where, writing \(u=1+y\),
\begin{align*}
 P(1+y)={}&16363+211070y+759380y^2+1274440y^3\\
           &+1128960y^4+515200y^5+96000y^6>0
           \qquad(y\geq0).
\end{align*}
Therefore \(F'_u(h)\leq b(u)-hK(u)<0\), while \(F'_u(0)>0\).
Strict concavity places the global maximum on
\([-\tfrac12,\tfrac12]\) inside \((0,h)\).

On \([0,h]\), integration of the second-derivative bound and
completion of the square yield
\[
 F_u(\delta)
 \leq F_u(0)+F'_u(0)\delta-\tfrac12K(u)\delta^2
 \leq F_u(0)+\frac{F'_u(0)^2}{2K(u)}.
\]
Lemma \ref{center} gives \(F_u(0)=c(u)-1/(8u)<H(u)-1/(8u)\).
Using \eqref{slope}, the global maximum is consequently less than
\(-R(u)\), where
\begin{equation}\label{R}
 R(u)=\frac1{8u}-H(u)-\frac{b(u)^2}{2K(u)}.
\end{equation}
It remains to check a rational inequality. Direct algebra gives
\begin{equation}\label{certificate}
 R(u)=
 \frac{N(u)}
 {129024000u^8(u+1)^5(2u+1)^4(300u^2+60u+53)}.
\end{equation}
The following table gives every coefficient \(c_j\) in
\(N(1+y)=\sum_{j=0}^{16}c_jy^j\):
\[
\begin{array}{r|r@{\qquad}r|r}
j&c_j&j&c_j\\ \hline
0&32048581152&9&124606427560816\\
1&673338924432&10&65045265578720\\
2&5218137436272&11&26078831078080\\
3&22582411828872&12&7883010992000\\
4&63881423270522&13&1738707609600\\
5&127854445158661&14&264182784000\\
6&189369613243226&15&24729600000\\
7&213134658410116&16&1075200000\\
8&184998793083656&&
\end{array}
\]
Every coefficient is positive. Hence \(R(u)>0\) for all \(u\geq1\),
which proves \(F_u(\delta)<0\) on the entire required interval and
establishes \eqref{target}.

To prove optimality, specialize to \(\alpha=1/2\), so
\(I(x,1/2)=c(u)\) and \(x+\alpha=u-1/2\).
For any fixed \(C>4\), Lemma \ref{center} gives
\[
 c(u)-\frac1{8u+C-4}
 =\frac{C-4}{64u^2}+O(u^{-3})>0
\]
for all sufficiently large \(u\). Such a constant \(C\) cannot
give a valid upper bound. The proved inequality with \(C=4\)
therefore has the largest possible additive constant.
\end{proof}

\noindent\textbf{Exact verification and scope.}
The companion \texttt{verify\_polynomials.py} uses only exact
standard-library integer and rational arithmetic. It constructs the
rational functions \eqref{J} and \eqref{R} directly from their
definitions, verifies both displayed polynomial identities by formal
cross-multiplication, and checks the shifted coefficients by independent
binomial and Horner substitutions. Both checks concern exact
polynomials, so they certify the infinite range without a numerical grid.
The main identity was also derived independently with symbolic
rational simplification.

The argument includes \(x=0\), \(\alpha=0\), \(\alpha=1\), and
arbitrarily large \(x\). All denominators are strictly positive.
The strict midpoint bound and the positive polynomial give strictness
at every finite point; sharpness of the additive constant arises at
infinity. No part of the source's conjectured inequality is assumed.

\section*{Tasks 2--5. Further directions after the full solve}
These are proposed tasks rather than additional solved assertions.
Check current literature and public claims before attacking each one.
\begin{description}
\item[Task 2.] Determine the infimum of the normalized deficit
\[
 \bigl(x+\alpha+\tfrac12\bigr)^3
 \left(\frac1{8(x+\alpha)+4}-I(x,\alpha)\right)
\]
on the full domain, and identify its minimizing points or limiting
sequences.
\item[Task 3.] For each \(x\geq0\), determine the unique maximizer
of the difference between \(I(x,\alpha)\) and its upper bound.
Prove monotonicity properties of this maximizer and obtain a
uniformly controlled asymptotic expansion.
\item[Task 4.] Classify the parameter pairs in the larger domain
\(x>-1,\ 0\leq\alpha\leq1,\ x+\alpha>-1/2\) for which the same
bound remains valid.
\item[Task 5.] Find a direct integral or monotonicity proof of the
sharp additive constant \(4\) that also proves the full bound
without the degree-six and degree-sixteen polynomial certificates.
\end{description}

\begin{thebibliography}{9}
\bibitem{Schmidlin}
M. Schmidlin, \emph{On Gautschi \& Stirling Identities, Asymptotics
and Inequalities for the Pi (or Gamma) Function},
arXiv:2601.21906v1 (2026),
\url{https://arxiv.org/abs/2601.21906}.
\bibitem{DLMF}
NIST Digital Library of Mathematical Functions,
Eqs.~5.9.10, 5.9.13 and 5.9.16,
\url{https://dlmf.nist.gov/5.9}.
\end{thebibliography}
\end{document}
